% Figure: bend-to-tie, held-out set 2, per model, with Wilson 95% intervals (paper/ci.py wilson()).
% Sources: paper/heldout2_v2.md per-run bend rows summed by model (draft.md Section 6 table):
%   own checks tie: Opus 398/423 (91.4-96.0), Haiku 454/1,539 (27.3-31.8), luna 272/2,805 (8.7-10.8)
%   erased: Opus 10/13 (49.7-91.8), Haiku 7/13 (29.1-76.8), luna 4/13 (12.7-57.6)
\begin{figure}[t]
\centering
\begin{tikzpicture}
\begin{groupplot}[group style={group size=2 by 1, horizontal sep=16pt, yticklabels at=edge left},
  dotplot, width=2.75cm, y=13pt, xmin=0, xmax=100,
  ytick={1,2,3}, ymin=1, ymax=3, y dir=reverse,
  yticklabels={Claude Opus, Claude Haiku, gpt-6-luna},
  xtick={0,50,100}, xticklabels={0,50,100\%}, enlarge y limits={abs=0.75}, title style={text height=1.6ex, text depth=0.4ex}]
\nextgroupplot[title={own checks tie}]
\addplot[dotmark, error bars/.cd, x dir=both, x explicit] coordinates {
  (94.09,1) -= (2.67,0) += (1.88,0)
  (29.50,2) -= (2.23,0) += (2.33,0)
  (9.70,3) -= (1.04,0) += (1.15,0)};
\coordinate (ben1) at (axis cs:94.09,1);
\coordinate (ben2) at (axis cs:29.50,2);
\coordinate (ben3) at (axis cs:9.70,3);
\nextgroupplot[title={planted issues erased}]
\addplot[dotmark, error bars/.cd, x dir=both, x explicit] coordinates {
  (76.92,1) -= (27.18,0) += (14.90,0)
  (53.85,2) -= (24.70,0) += (22.95,0)
  (30.77,3) -= (18.09,0) += (26.86,0)};
\coordinate (ben4) at (axis cs:76.92,1);
\coordinate (ben5) at (axis cs:53.85,2);
\coordinate (ben6) at (axis cs:30.77,3);
\end{groupplot}
\node[vlab] at (ben1) {398/423};
\node[vlab] at (ben2) {454/1,539};
\node[vlab] at (ben3) {272/2,805};
\node[vlab] at (ben4) {10/13};
\node[vlab] at (ben5) {7/13};
\node[vlab] at (ben6) {4/13};
\end{tikzpicture}
\caption{Bending to tie on the private-equity held-out set~2 (bad scans, five funds per model), with Wilson 95\% intervals. Left: share of tier~0 errors whose checks all pass on the model's own numbers. Right: share of the 13 planted document issues per model that the model's entries erased.}
\label{fig:bend}
\end{figure}
